% TOP_PROBLEM: {"id": 30006765, "title": "Capacity of the General Discrete Memoryless Relay Channel", "category_id": 15, "category": {"id": 15, "name": "computer_science", "display_name": "Computer Science", "description": "Computational complexity, algorithms, and theoretical CS.", "slug": "computer-science", "order_index": 15, "created_at": "2026-07-31T15:26:25.671Z"}, "status": "open"}
% FIRST_POSTED: 2026-09-25
% ENTRY_KIND: statement_only
% !TeX program = lualatex
% Definition and sources only; this file is not an LLM solution attempt.
\documentclass[11pt]{article}
\usepackage[margin=25mm]{geometry}
\usepackage{fontspec}
\setmainfont{Latin Modern Roman}
\usepackage{amsmath,amssymb,mathtools}
\usepackage{unicode-math}
\setmathfont{Latin Modern Math}
\usepackage{xurl,hyperref}
\hypersetup{colorlinks=true,urlcolor=blue}
\setcounter{secnumdepth}{0}
\setlength{\parindent}{0pt}
\setlength{\parskip}{5pt}
\setlength{\emergencystretch}{3em}
\begin{document}
\section*{\texorpdfstring{Capacity of the General Discrete Memoryless Relay Channel}{Capacity of the General Discrete Memoryless Relay Channel}}\label{30006765}
\subsection{Definitions and mathematical statement}
A discrete memoryless relay channel has finite alphabets and a conditional law $W(y,y_r\mid x,x_r)$. A length-$n$ code consists of a source encoder for a uniform message $M$, causal relay encoders $x_{r,i}=f_i(y_r^{i-1})$, and a destination decoder $\widehat M=g(y^n)$. A rate $R$ is achievable when codes exist with $\liminf n^{-1}\log_2|\mathcal M_n|\ge R$ and $\Pr(\widehat M\ne M)\to0$. Define
\[
 C(W)=\sup\{R:R\text{ is achievable}
 \}.
\]
The problem asks for a computable exact characterization for every $W$, not this operational definition repeated as an answer. The relay cannot depend on its current or future observation. Average error, not zero error, is the criterion.
\par\smallskip\textit{Formulation and concept references:} \href{https://arxiv.org/html/2405.09534v3}{[S1]}.

\subsection{Short English statement}
What is the exact reliable communication rate for an arbitrary finite relay channel when the relay can use only its past observations?
\subsection{Sources}
\begin{itemize}
\item[S1]Ozyilkan, Carpi, Garg and Erkip, Learning-Based Compress-and-Forward Schemes for the Relay Channel.
\url{https://arxiv.org/html/2405.09534v3}.
\textit{Formulation source.}
\item[S2]Status source for Capacity of the General Discrete Memoryless Relay Channel.
\url{https://arxiv.org/html/2609.15709v1}.
\textit{Further reference; not a proof endorsement.}
\end{itemize}
\end{document}
